% 3.5 230 个三维空间群的分类
\section{三维空间群 230 种的分类}

到目前为止，我们在本章中讨论的都是空间群的平移对称性。空间群 $\mathbf{G}$ 有一个三维的纯平移子群 $\mathbf{T}$，它确定布拉维格子。布拉维格子具有 $\mathbf{T}$ 的对称性；此外它还在点群 $\mathbf{P}$ 的操作下不变。事实上，若 $\mathbf{F}$ 是布拉维格子的全对称群，则 $\mathbf{F} = \mathbf{T} \wedge \mathbf{P}$，即 $\mathbf{T}$ 与 $\mathbf{P}$ 的半直积（半直积的含义见定义 (1.2.17)）。$\mathbf{P}$ 确定 $\mathbf{G}$ 的晶系。$\mathbf{F}$ 本身是一个空间群。若 $\mathbf{G}$ 是点式的，则 $\mathbf{G}$ 要么与 $\mathbf{F}$ 重合，要么是 $\mathbf{F}$ 的子群；这就给出 73 个空间群，见 1.5 节。显然此时 $\mathbf{G} = \mathbf{T} \wedge \mathbf{Q}$，其中 $\mathbf{Q}$ 是 $\mathbf{P}$ 所属晶系中的一个点群。$\mathbf{Q}$ 确定该晶体的点群，而空间群为 $\mathbf{G}$ 的晶体就属于这个点群。

其余 157 个非点式空间群可以这样得到：把 $\mathbf{Q}$ 的一部分元素换成螺旋轴转动或滑移面反射。此时 $\mathbf{Q}$ 不再是群，也就不能再写 $\mathbf{G} = \mathbf{T} \wedge \mathbf{Q}$；这是因为，若集合 $\mathbf{Q}$ 的元素能够取成一组构成群，则它将同构于某个点群，从而 $\mathbf{G}$ 又将同构于某个点式空间群。

现在考虑必须对 $\mathbf{Q}$ 的元素 $\{R \mid \mathbf{v}\}$ 施加哪些限制，才能使 $\mathbf{G}$ 成为一个空间群。以下讨论对点式空间群显然成立，因此可以视为完整的刻画。首先 $\mathbf{G}$ 必须是群。于是，若 $\{R \mid \mathbf{v}\}$ 是 $\mathbf{G}$ 的元素，则其逆 $\{R^{-1} \mid -R^{-1}\mathbf{v}\}$ 也必须是 $\mathbf{G}$ 的元素（见式 (1.5.3)）。因此，若 $\{R \mid \mathbf{w}\}$ 也在 $\mathbf{G}$ 中，则
\begin{equation*}
\{R^{-1} \mid -R^{-1}\mathbf{v}\}\{R \mid \mathbf{w}\} = \{E \mid R^{-1}\mathbf{w} - R^{-1}\mathbf{v}\}
\tag{3.5.1}
\end{equation*}
也是 $\mathbf{G}$ 的元素。事实上它必属于平移子群 $\mathbf{T}$，因为其转动部分是恒等元。其次，$\mathbf{G}$ 有一个属于某一晶系的布拉维格子，该晶系的全对称点群 $\mathbf{P}$ 必须包含操作 $R$。于是，如式 (3.5.1) 所蕴含的，若 $R^{-1}(\mathbf{w} - \mathbf{v})$ 是 $\mathbf{F}$ 的一个平移，则 $(\mathbf{w} - \mathbf{v})$ 也是。这意味着与任一点群操作相联系的平移部分相差的只能是 $\mathbf{T}$ 的元素。因此可以用 $\mathbf{T}$ 的左陪集代表把 $\mathbf{G}$ 写成（见定义 1.2.9）
\begin{equation*}
\mathbf{G} = \{R_{1} \mid \mathbf{v}_{1}\}\mathbf{T} + \{R_{2} \mid \mathbf{v}_{2}\}\mathbf{T} + \dots + \{R_{h} \mid \mathbf{v}_{h}\}\mathbf{T},
\tag{3.5.2}
\end{equation*}
并且可以选择 $R_{1} = E$，而 $\mathbf{v}_{i}$ 取为与 $R_{i}$（$i = 1$ 至 $h$）相联系的、模长尽可能小的矢量。于是对点式空间群所有的 $\mathbf{v}_{i}$ 都是零，而对非点式空间群则必须有一个或多个 $\mathbf{v}_{i}$ 非零。$\mathbf{Q}$ 由陪集代表 $\{R_{i} \mid \mathbf{v}_{i}\}$ 组成。$h$ 个元素 $R_{1}, R_{2}, \dots, R_{h}$ 构成一个点群 $\mathbf{F}$，即 $\mathbf{G}$ 的同形点群（见 1.5 节）。$\mathbf{F}$ 是 $\mathbf{P}$ 的子群。对点式空间群，$\mathbf{F}$ 与 $\mathbf{Q}$ 重合。$h$ 称为 $\mathbf{G}$ 的\textit{宏观阶}（macroscopic order），它是 $\mathbf{T}$ 在 $\mathbf{G}$ 中的指数。

在上述描述中，非零矢量 $\mathbf{v}$ 并非任意的，因为要满足 $\mathbf{G}$ 的群的要求。例如，若 $R$ 的阶为 $n$，则 $\{R \mid \mathbf{v}\}^{n}$ 必须属于 $\mathbf{T}$，即 $(\mathbf{v} + R\mathbf{v} + \dots + R^{n-1}\mathbf{v})$ 必须是一个平移。这一要求以及其他类似要求把非点式空间群的数目限制为 157 个，总计 230 个空间群。

\noindent\textbf{定理 3.5.1}\quad $\mathbf{T}$ 是 $\mathbf{G}$ 的不变子群。因为若 $\{R \mid \mathbf{v}\}$ 是 $\mathbf{G}$ 的元素，$\{E \mid \mathbf{t}\}$ 是 $\mathbf{T}$ 的元素，则
\begin{equation*}
\{R \mid \mathbf{v}\}\{E \mid \mathbf{t}\}\{R \mid \mathbf{v}\}^{-1} = \{R \mid \mathbf{v} + R\mathbf{t}\}\{R^{-1} \mid -R^{-1}\mathbf{v}\} = \{E \mid R\mathbf{t}\}
\tag{3.5.3}
\end{equation*}
属于 $\mathbf{T}$，因为 $R\mathbf{t}$ 是 $\mathbf{T}$ 的一个平移。

每个空间群诸元素的 Seitz 空间群符号已由许多作者给出（Faddeyev 1964, Koptsik 1966, Kovalev 1965, Lyubarskii 1960, Miller and Love 1967），也可以利用《X 射线晶体学国际表》（Henry and Lonsdale 1965）写出来。在该表的第 4.3 节中，对每个空间群给出了晶胞中等价位置的一组坐标；对每个非立方空间群还给出了图示其对称元素的图。利用这些表不仅能够得到以 Seitz 符号 $\{R_{i} \mid \mathbf{v}_{i}\}$ 表示的空间群元素，还能对这些对称操作作出几何解释；例如对一根螺旋转动轴，可以确定轴的方向以及转动的大小和沿轴平移的大小（Wondratschek and Neubüser 1967）。我们将通过例子说明《国际表》中等价位置的这两种用法。

我们以基于正交初基布拉维格子的空间群 $Pmmn$（$D_{2h}^{13}$）为例，说明如何利用《国际表》给出的等价位置确定相应的 Seitz 空间群符号。等价位置既对晶胞中完全一般的点给出，也对位于各种对称元素上的点给出。图 3.17 复制了《国际表》第一卷第 4.3 节中与 $Pmmn$（$D_{2h}^{13}$）有关的两页。在图 3.17(a) 中，第一组等价位置对应于晶胞中的一般点 $(x, y, z)$；其余各组等价位置对应于位于某个对称轴或对称面上的点 $(x, y, z)$，这些与我们无关。从点 $\mathbf{r}_{1}\,(=(x, y, z))$ 出发，施加由 $\{R_{1} \mid \mathbf{v}_{1}\}$ 表示的空间群操作，由式 (1.5.2) 得
\begin{equation*}
\mathbf{r}_{1}' = R_{1}\mathbf{r}_{1} + \mathbf{v}_{1}
\tag{3.5.4}
\end{equation*}
所得的矢量 $\mathbf{r}_{1}'$ 为
\begin{center}
\begin{tabular}{cccc}
$x, y, z$;\quad & $\bar{x}, \bar{y}, z$;\quad & $\bar{x}, y, z$;\quad & $x, \bar{y}, z$;\quad \\[0.5ex]
$\frac{1}{2}-x, \frac{1}{2}-y, \bar{z}$;\quad & $\frac{1}{2}-x, \frac{1}{2}+y, \bar{z}$;\quad & $\frac{1}{2}+x, \frac{1}{2}+y, \bar{z}$;\quad & $\frac{1}{2}+x, \frac{1}{2}-y, \bar{z}$.
\end{tabular}
\end{center}
这些分数是指《国际表》中对每个空间群（立方空间群除外）所画的惯用晶胞在 $x$、$y$、$z$ 方向的棱长的分数。利用表 1.4 可以对这八个 $\mathbf{r}_{1}'$ 逐一辨认出 $R_{1}$：
\begin{center}
\begin{tabular}{cccccccc}
$xyz$;\quad & $\bar{x}\bar{y}z$;\quad & $\bar{x}yz$;\quad & $x\bar{y}z$;\quad & $\bar{x}\bar{y}\bar{z}$;\quad & $\bar{x}y\bar{z}$;\quad & $xy\bar{z}$;\quad & $x\bar{y}\bar{z}$;\\
$E$;\quad & $C_{2z}$;\quad & $\sigma_{x}$;\quad & $\sigma_{y}$;\quad & $I$;\quad & $C_{2y}$;\quad & $\sigma_{z}$;\quad & $C_{2x}$
\end{tabular}
\end{center}
这就给出了空间群符号中的转动部分。与其中前四个相联系的平移矢量 $\mathbf{v}_{1}$ 显然为零。由于对正交初基布拉维格子有 $\mathbf{t}_{1} = (0, -b, 0)$、$\mathbf{t}_{2} = (a, 0, 0)$（见表 3.1），与后四个符号相联系的平移矢量是 $-\frac{1}{2}\mathbf{t}_{1} + \frac{1}{2}\mathbf{t}_{2}$。可以在平移矢量 $\mathbf{v}_{1}$ 上加上 $\mathbf{t}_{1}$、$\mathbf{t}_{2}$ 或 $\mathbf{t}_{3}$ 的整数倍，因此我们取 $\mathbf{v}_{1} = +\frac{1}{2}\mathbf{t}_{1} + \frac{1}{2}\mathbf{t}_{2}$，于是空间群 $Pmmn$（$D_{2h}^{13}$）的元素为
\begin{center}
$\{E \mid 000\}$，$\{C_{2z} \mid 000\}$，$\{\sigma_{x} \mid 000\}$，$\{\sigma_{y} \mid 000\}$，\\[0.5ex]
$\{I \mid \frac{1}{2}\frac{1}{2}0\}$，$\{C_{2y} \mid \frac{1}{2}\frac{1}{2}0\}$，$\{\sigma_{z} \mid \frac{1}{2}\frac{1}{2}0\}$，以及 $\{C_{2x} \mid \frac{1}{2}\frac{1}{2}0\}$。
\end{center}
不必把所有这些元素都列出来：在表 3.7 中我们只列出了每个空间群的生成元。在 $Pmmn$（$D_{2h}^{13}$）名下可以看到，选取的生成元是 $\{C_{2x} \mid \frac{1}{2}\frac{1}{2}0\}$、$\{C_{2y} \mid \frac{1}{2}\frac{1}{2}0\}$ 和 $\{I \mid \frac{1}{2}\frac{1}{2}0\}$。

\begin{figure}[H]
\centering
\includegraphics[width=0.72\textwidth]{fig3_17a_mu.jpg}
\caption*{\textbf{图 3.17(a)}\ 《国际表》（Henry and Lonsdale 1965）书影：$Pmmn$（$D_{2h}^{13}$）的第一种定向（原点在 $mmn$）。}
\end{figure}
\begin{figure}[H]
\centering
\includegraphics[width=0.72\textwidth]{fig3_17b_mu.jpg}
\caption*{\textbf{图 3.17(b)}\ 同上，第二种定向（原点在 $\bar{1}$）。}
\end{figure}

重要的是要弄清：若把描述晶体所用的坐标系原点移到晶胞中另一点，给定的 Seitz 空间群符号 $\{R_{1} \mid \mathbf{v}_{1}\}$ 的形式会受到什么影响。设 $\mathbf{r}_{1}$ 与 $\mathbf{r}_{1}'$ 是某个点在施加一个主动空间群对称操作前、后的位置；在第一组坐标轴 $O_{1}x_{1}y_{1}z_{1}$ 中该操作记为 $\{R_{1} \mid \mathbf{v}_{1}\}$。现在另取一组坐标轴 $O_{2}x_{2}y_{2}z_{2}$，其原点相对 $O_{1}x_{1}y_{1}z_{1}$ 的原点移动了矢量 $+\mathbf{t}_{0}$ 而轴不转动，见图 3.18。由该图可见
\begin{equation*}
\mathbf{r}_{1} = \mathbf{t}_{0} + \mathbf{r}_{2}
\tag{3.5.5}
\end{equation*}
及
\begin{equation*}
\mathbf{r}_{1}' = \mathbf{t}_{0} + \mathbf{r}_{2}',
\tag{3.5.6}
\end{equation*}
又由在 $O_{1}x_{1}y_{1}z_{1}$ 中使用的空间群符号 $\{R_{1} \mid \mathbf{v}_{1}\}$ 的定义，
\begin{equation*}
\mathbf{r}_{1}' = R_{1}\mathbf{r}_{1} + \mathbf{v}_{1}.
\tag{3.5.4}
\end{equation*}
若这一空间群操作在坐标系 $O_{2}x_{2}y_{2}z_{2}$ 中记为 $\{R_{2} \mid \mathbf{v}_{2}\}$，则 $\mathbf{r}_{2}'$ 与 $\mathbf{r}_{2}$ 满足
\begin{equation*}
\mathbf{r}_{2}' = R_{2}\mathbf{r}_{2} + \mathbf{v}_{2}
\tag{3.5.7}
\end{equation*}
于是可以用这些方程把 $\{R_{2} \mid \mathbf{v}_{2}\}$ 用 $\{R_{1} \mid \mathbf{v}_{1}\}$ 表示出来。由式 (3.5.7)，
\begin{align*}
R_{2}\mathbf{r}_{2} + \mathbf{v}_{2} &= \mathbf{r}_{2}' = \mathbf{r}_{1}' - \mathbf{t}_{0} \quad\text{（由 (3.5.6)）}\\
&= R_{1}\mathbf{r}_{1} + \mathbf{v}_{1} - \mathbf{t}_{0} \quad\text{（由 (3.5.4)）}\\
&= R_{1}(\mathbf{t}_{0} + \mathbf{r}_{2}) + \mathbf{v}_{1} - \mathbf{t}_{0} \quad\text{（由 (3.5.5)）}\\
&= R_{1}\mathbf{r}_{2} + \mathbf{v}_{1} + R_{1}\mathbf{t}_{0} - \mathbf{t}_{0},
\end{align*}
即
\begin{equation*}
R_{2}\mathbf{r}_{2} + \mathbf{v}_{2} = R_{1}\mathbf{r}_{2} + \mathbf{v}_{1} + R_{1}\mathbf{t}_{0} - \mathbf{t}_{0}.
\tag{3.5.8}
\end{equation*}
比较式 (3.5.8) 两边可见，把原点移动 $+\mathbf{t}_{0}$ 对空间群符号 $\{R_{1} \mid \mathbf{v}_{1}\}$ 的效果是产生 $\{R_{2} \mid \mathbf{v}_{2}\}$，其中
\begin{equation*}
R_{2} = R_{1}
\tag{3.5.9}
\end{equation*}
及
\begin{equation*}
\mathbf{v}_{2} = \mathbf{v}_{1} + R_{1}\mathbf{t}_{0} - \mathbf{t}_{0}.
\tag{3.5.10}
\end{equation*}
也就是说，移动原点 $+\mathbf{t}_{0}$ 对空间群符号 $\{R_{1} \mid \mathbf{v}_{1}\}$ 的效果是产生
\begin{equation*}
\{R_{2} \mid \mathbf{v}_{2}\} = \{R_{1} \mid \mathbf{v}_{1} + R_{1}\mathbf{t}_{0} - \mathbf{t}_{0}\}.
\tag{3.5.11}
\end{equation*}

\begin{figure}[H]
\centering
\includegraphics[width=0.55\textwidth]{fig3_18_mu.jpg}
\caption*{\textbf{图 3.18}\ 通过一个平移改变原点。}
\end{figure}

我们仍以空间群 $Pmmn$（$D_{2h}^{13}$）为例说明移动原点对空间群符号的影响。《国际表》对该空间群给出了两种不同的定向，见图 3.17。上面我们已用第一种定向（图 3.17(a)）写出了空间群元素。在第二种定向（图 3.17(b)）中，原点移动了（沿 $x$ 轴重复距离的 $-\frac{1}{4}$）加（沿 $y$ 轴重复距离的 $-\frac{1}{4}$），即由表 3.1，
\begin{equation*}
\mathbf{t}_{0} = \frac{1}{4}\mathbf{t}_{1} - \frac{1}{4}\mathbf{t}_{2}.
\tag{3.5.12}
\end{equation*}
于是利用式 (3.5.9) 与 (3.5.10)，可以把先前得到的空间群元素转换到第二组轴上。例如 $\{C_{2z} \mid 000\}$ 被 $\{R_{2} \mid \mathbf{v}_{2}\}$ 取代，其中由式 (3.5.9)，$R_{2} = C_{2z}$，而
\begin{align*}
\mathbf{v}_{2} &= \mathbf{v}_{1} + R_{1}\mathbf{t}_{0} - \mathbf{t}_{0}\\
&= \mathbf{0} + C_{2z}\left(\tfrac{1}{4}\mathbf{t}_{1} - \tfrac{1}{4}\mathbf{t}_{2}\right) - \left(\tfrac{1}{4}\mathbf{t}_{1} - \tfrac{1}{4}\mathbf{t}_{2}\right)\\
&= \left(-\tfrac{1}{4}\mathbf{t}_{1} + \tfrac{1}{4}\mathbf{t}_{2}\right) - \left(\tfrac{1}{4}\mathbf{t}_{1} - \tfrac{1}{4}\mathbf{t}_{2}\right)\quad\text{（由表 3.2）}\\
&= -\tfrac{1}{2}\mathbf{t}_{1} + \tfrac{1}{2}\mathbf{t}_{2},
\tag{3.5.13}
\end{align*}
故新坐标系中的元素为 $\{C_{2z} \mid \bar{\frac{1}{2}}\bar{\frac{1}{2}}0\}$，可以加上 $\mathbf{t}_{1}$ 得到以新原点为基准的空间群元素 $\{C_{2z} \mid \frac{1}{2}\frac{1}{2}0\}$。这与图 3.17(b) 所示的等价位置一致：$\{C_{2z} \mid \frac{1}{2}\frac{1}{2}0\}$ 作用在 $(x, y, z)$ 上产生 $(\frac{1}{2}-x, \frac{1}{2}-y, z)$，这正是图 3.17(b) 给出的八个一般等价位置之一。对同样的 $\mathbf{t}_{0} = \frac{1}{4}\mathbf{t}_{1} - \frac{1}{4}\mathbf{t}_{2}$ 反复使用式 (3.5.9) 与 (3.5.10)，得到以新原点为基准的空间群 $Pmmn$（$D_{2h}^{13}$）的全部元素：
\begin{center}
$\{E \mid 000\}$，$\{C_{2z} \mid \frac{1}{2}\frac{1}{2}0\}$，$\{\sigma_{x} \mid 0\frac{1}{2}0\}$，$\{\sigma_{y} \mid \frac{1}{2}00\}$，\\[0.5ex]
$\{I \mid 000\}$，$\{\sigma_{z} \mid \frac{1}{2}\frac{1}{2}0\}$，$\{C_{2x} \mid 0\frac{1}{2}0\}$，$\{C_{2y} \mid \frac{1}{2}00\}$。
\end{center}
在这个清单里，我们在某些情况下加上或减去了正交初基布拉维格子平移群的矢量，以简化空间群元素的形式。

230 个空间群中每一个的生成元都列于表 3.7。在前三列中我们用国际符号和 Schönflies 符号标识空间群，第 4 列给出足以确定空间群 $\mathbf{G}$ 全部元素的 Seitz 空间群符号，其中 $\{R \mid pqr\}$ 表示 $\{R \mid p\mathbf{t}_{1} + q\mathbf{t}_{2} + r\mathbf{t}_{3}\}$，$\mathbf{t}_{1}$、$\mathbf{t}_{2}$、$\mathbf{t}_{3}$ 按表 3.1 的定义。这些操作一如既往地用 1.5 节解释的 Seitz 记号表示（见式 (1.5.2) 与 (1.5.3)）。

\begin{figure}[H]
\centering
\includegraphics[width=0.92\textwidth]{c3_t37_p1.png}
\caption*{\textbf{表 3.7}\ 230 个空间群（原书第 126 页起扫描；第 1 页）。}
\end{figure}
\begin{figure}[H]
\centering
\includegraphics[width=0.92\textwidth]{c3_t37_p2.png}
\caption*{\textbf{表 3.7}（第 2 页）。}
\end{figure}
\begin{figure}[H]
\centering
\includegraphics[width=0.92\textwidth]{c3_t37_p3.png}
\caption*{\textbf{表 3.7}（第 3 页）。}
\end{figure}
\begin{figure}[H]
\centering
\includegraphics[width=0.92\textwidth]{c3_t37_p4.png}
\caption*{\textbf{表 3.7}（第 4 页）。}
\end{figure}
\begin{figure}[H]
\centering
\includegraphics[width=0.92\textwidth]{c3_t37_p5.png}
\caption*{\textbf{表 3.7}（第 5 页）。}
\end{figure}
\begin{figure}[H]
\centering
\includegraphics[width=0.92\textwidth]{c3_t37_p6.png}
\caption*{\textbf{表 3.7}（第 6 页）。}
\end{figure}
\begin{figure}[H]
\centering
\includegraphics[width=0.92\textwidth]{c3_t37_p7.png}
\caption*{\textbf{表 3.7}（第 7 页）。}
\end{figure}
\begin{figure}[H]
\centering
\includegraphics[width=0.92\textwidth]{c3_t37_p8.png}
\caption*{\textbf{表 3.7}（第 8 页，页末起为表的注）。}
\end{figure}
\begin{figure}[H]
\centering
\includegraphics[width=0.92\textwidth]{c3_t37_p9.png}
\caption*{\textbf{表 3.7}（注 (ii) 续--(vi) 及取向符号表，原书扫描）。}
\end{figure}
\begin{figure}[H]
\centering
\includegraphics[width=0.80\textwidth]{c3_t37_p10.png}
\caption*{\textbf{表 3.7}（注 (vi) 附表：受 Buerger (1967) 建议影响的空间群，原书扫描）。}
\end{figure}

表 3.7 的注（译文）：

\noindent（i）我们在第 1、2、3 列中分别给出每个空间群的国际编号、国际符号和 Schönflies 符号。第 3 列的 Schönflies 符号前冠以相应布拉维格子的符号（见表 3.1）。

\noindent（ii）第 1、2、3 列给出的空间群名称蕴含（a）它所依托的布拉维格子，以及（b）同形点群。因此无需单独列出平移群 $\mathbf{T}$ 的元素（见表 3.1）或同形点群 $\mathbf{F}$ 的元素（见表 2.2）。尚需给出的是若干陪集代表（即 3.5 节所述 $\mathbf{Q}$ 的元素）。我们给出的数目足以确定空间群 $\mathbf{G}$ 的全部元素；这些元素列于第 4 列，其中 $\{R \mid pqr\}$ 表示 $\{R \mid p\mathbf{t}_{1} + q\mathbf{t}_{2} + r\mathbf{t}_{3}\}$，$\mathbf{t}_{1}$、$\mathbf{t}_{2}$、$\mathbf{t}_{3}$ 按表 3.1 的定义。这些操作一如既往地用 1.5 节解释的 Seitz 记号表示（见式 (1.5.2) 与 (1.5.3)）。

\noindent（iii）对某些空间群，我们所用的元素与《国际表》（Henry and Lonsdale 1965）所给的略有不同。对这些情形，我们在第 4 列的第一行给出我们用于 $\mathbf{Q}$ 的元素的 Seitz 空间群符号，第二行给出《国际表》所用的元素。

\noindent（iv）尽管我们保留编号 38--41（$Amm2\ (C_{2v}^{14})$、$Abm2\ (C_{2v}^{15})$、$Ama2\ (C_{2v}^{16})$、$Aba2\ (C_{2v}^{17})$）的国际记号，但不使用 A 格子，因为那将需要在表 3.6 的 $\Gamma_{o}^{b}$ 之下增加一个条目，并在图 3.6 中增加一幅图（见表 3.1 的注 (v)）。取而代之，我们对所有单斜底心空间群一律使用 C 格子，并把编号 38--41 的点群操作的取向改为使二次转动轴沿 $y$ 轴，而在《国际表》中它沿 $z$ 轴。

\noindent（v）若第 5 列给出一个平移 $\mathbf{t}_{0}$，则把它代入式 (3.5.11) 就能把我们所用的元素 $\{R_{1} \mid \mathbf{v}_{1}\}$ 变换为《国际表》的元素 $\{R_{2} \mid \mathbf{v}_{2}\}$。第 5 列中的星号表示我们的原点与《国际表》的原点重合，但轴的取向不同。若第 5 列既有矢量 $\mathbf{t}_{0}$ 又有星号，则表示矢量 $\mathbf{t}_{0}$ 按式 (3.5.11) 把我们所用的 $\{R_{1} \mid \mathbf{v}_{1}\}$ 变换为 $\{R_{2} \mid \mathbf{v}_{2}\}$，后者与《国际表》的元素只在取向上不同。

\noindent（vi）Buerger（1967）对某些空间群的国际符号提出了修改。他呼吁回到旧版《国际表》的做法：晶格与对称轴的相对取向完全由空间群符号中的布拉维格子部分表明，而完全不通过符号中的点群（或其替代）部分体现。若用符号 E 表示端心格子而不蕴含取向，则正交空间群中 E 的可取取向为 $A$、$B$ 和 $C$。若用 S 和 Z 表示简单和体心格子而不蕴含取向，则表明取向的相应符号如第 9 页扫描图中的表所示。对四方空间群，需要对 $S$ 的两种可取取向（P 和 D）以及 $Z$ 的两种可能取向（I 和 J）使用替代符号；对三方晶体，需要对 $S$ 的三种可取取向（P、D 和 R）使用替代符号。Buerger 敦促采用这些格子符号，并取缔符号“1”的使用——点群 $1\ (C_{1})$ 和 $\bar{1}\ (C_{i})$ 以及空间群 $P1\ (C_{1}^{1})$ 和 $P\bar{1}\ (C_{i}^{1})$ 除外。只有注 (vi) 附表中列出的空间群会受到这些改动的影响。经 Buerger（1967）引入这些修改后，任何空间群标签的第一个符号都表明布拉维格子的一般类型及其相对对称轴的取向，标签中其余的符号则要么是点群符号，要么是其替代符号。这些修改的一个结果是：对每个空间群，这些符号的顺序都将与相应点群的标准顺序完全一致。

空间群本身虽由《国际表》正式定义，但原点的选取和坐标轴取向的选择仍留有若干自由；《国际表》本身对许多空间群就给出了两种可替换的定向，正如 $Pmmn$（$D_{2h}^{13}$）的例子所示。援引《国际表》第 52 页的话：“所采用的空间群标准定向并不具有官方的重要性。”对某些空间群，《新国际表》（Henry and Lonsdale 1965）所用的空间群符号曾受到 Buerger（1967）的批评。Buerger 建议的修改已列于表 3.7 的注 (vi)；若被采纳，将只影响少数空间群。

国际空间群符号的逻辑在 1.5 节已有简要叙述。这些符号用于《国际表》第一版（Bragg, von Laue, and Hermann 1935\textit{a}, \textit{b}）并逐渐通用。当时符号的布拉维格子部分并不完全令人满意，因此在编纂新《国际表》时作了一些修改（Henry and Lonsdale 1965）；表 3.7 采用的就是新版符号。例如，四方点群 $\bar{4}2m$ 放进初基四方布拉维格子的晶胞时，纯转动轴既可以平行于 $x$ 轴，也可以平行于直线 $x = y$。旧《国际表》用两种不同的布拉维格子符号 P 和 C 区分这两种取向；对体心四方格子上的两种取向则用 I 和 F 两种符号。新《国际表》中，同一个布拉维格子在其一切空间群中都用同一符号表示（编号 38--41 的空间群 $Amm2\ (C_{2v}^{14})$、$Abm2\ (C_{2v}^{15})$、$Ama2\ (C_{2v}^{16})$、$Aba2\ (C_{2v}^{17})$ 除外）。点群在布拉维格子上不同取向的区分问题，于是或者通过调换点群符号 $\bar{4}2m$ 和 $\bar{6}2m$ 中最后两个符号的次序（得到 $\bar{4}m2$ 和 $\bar{6}m2$）来解决，或者对三方空间群通过把符号“1”加在空间群符号的不同位置来解决。其结果是，对称性符号如今常常被部分地理解为不仅表明有哪些对称元素，还表明其取向，见 Henry and Lonsdale（1965）的表 3.3.2。国际点群标记各部分的身份问题在第七章还会涉及。

为了说明如何用《国际表》给出的等价位置确定对称元素的几何描述，我们遵循 Wondratschek and Neubüser（1967）的方法。应当从 1.4 与 1.5 节清楚地区分\textit{对称元素}（如转动轴或对称面，粗略地说这是晶体的一种物理属性）与\textit{对称操作}（它是晶体的覆盖操作，是点群或空间群这一抽象数学群的元素）。空间群对称操作 $\{R_{i} \mid \mathbf{r}_{i}\}$ 作用在点 $x, y, z$ 上，产生等价位置中的某一点，后者一般可写为
\begin{equation*}
\begin{aligned}
x', y', z' &= r_{1} + a_{11}x + a_{12}y + a_{13}z,\\
&\quad r_{2} + a_{21}x + a_{22}y + a_{23}z,\ r_{3} + a_{31}x + a_{32}y + a_{33}z,
\end{aligned}
\tag{3.5.14}
\end{equation*}
其中 $r_{i}$ 是有理数，$a_{ij}$ 为 $+1$、$-1$ 或 0。

\noindent\textbf{例 3.5.1}\quad 在 $Pmmn$（$D_{2h}^{13}$）的诸等价位置中有 $\frac{1}{2}+x, \frac{1}{2}-y, \bar{z}$；这意味着 $a_{11} = +1$，$a_{22} = a_{33} = -1$，$r_{1} = r_{2} = \frac{1}{2}$，$r_{3} = 0$，其余 $a_{ij} = 0$。

式 (3.5.14) 可以方便地写成矩阵形式
\begin{equation*}
\mathbf{M}\left(\{R_{i} \mid \mathbf{r}_{i}\}\right) = \begin{pmatrix} a_{11} & a_{12} & a_{13} & r_{1}\\ a_{21} & a_{22} & a_{23} & r_{2}\\ a_{31} & a_{32} & a_{33} & r_{3}\\ 0 & 0 & 0 & 1 \end{pmatrix} = (\mathbf{A}, \mathbf{r}_{i})
\tag{3.5.15}
\end{equation*}
其中 $\mathbf{M}(\{R_{i} \mid \mathbf{r}_{i}\})$ 作用在列矢量 $(x, y, z, 1)^{\mathrm{T}}$ 上。显然矩阵 $a_{ij}$ 已经以缩并的形式出现在表 1.4 的 Jones 符号中。$\{R_{i} \mid \mathbf{r}_{i}\}$ 的 $R_{i}$ 部分的几何意义当然可以直接由表 1.4 辨认；另一种做法是，$\mathbf{A}$ 的特征标 $\chi(\mathbf{A})$ 与行列式 $\det(\mathbf{A})$ 很容易计算，按表 3.8，这两个量完全刻画了 $R_{i}$ 的性质（Bethe 1929），其中 $C_{n}$ 表示转角为 $2\pi/n$ 的转动。

\begin{table}[H]
\centering
\caption*{\textbf{表 3.8}\ 由 $\chi(\mathbf{A})$ 和 $\det(\mathbf{A})$ 确定 $R_{i}$ 的形式}
\begin{tabular}{@{}lcccccccccc@{}}
\toprule
$R_{i}$ & $C_{1}$ & $C_{2}$ & $C_{3}$ & $C_{4}$ & $C_{6}$ & $IC_{1}$ & $IC_{2}$ & $IC_{3}$ & $IC_{4}$ & $IC_{6}$\\
 & & & & & & $\equiv I$ & $\equiv\sigma$ & $\equiv S_{6}$ & $\equiv S_{4}$ & $\equiv S_{3}$\\
\midrule
$\chi(\mathbf{A})$ & $3$ & $-1$ & $0$ & $1$ & $2$ & $-3$ & $1$ & $0$ & $-1$ & $-2$\\
$\det(\mathbf{A})$ & $+1$ & $+1$ & $+1$ & $+1$ & $+1$ & $-1$ & $-1$ & $-1$ & $-1$ & $-1$\\
\bottomrule
\end{tabular}
\end{table}

若 $\det(\mathbf{A}) = +1$，则对称元素或者是纯转动轴，或者是螺旋转动轴。于是我们要确定该轴的方向；对螺旋转动轴，还要确定螺旋矢量 $\mathbf{w}$，即平移 $\mathbf{r}_{i}$ 沿轴方向的分量。若 $\mathbf{q}$ 是过原点且平行于螺旋轴的矢量，则 $\mathbf{A}\mathbf{q} = \mathbf{q}$，因此给出螺旋轴方向的 $\mathbf{q}$ 可由
\begin{equation*}
(\mathbf{A} - \mathbf{E})\mathbf{q} = 0
\tag{3.5.16}
\end{equation*}
求得。对螺旋转动，为确定 $\mathbf{w}$，注意
\begin{equation*}
\{R_{i} \mid \mathbf{r}_{i}\}\mathbf{q} = \mathbf{A}\mathbf{q} + \mathbf{r}_{i} = \mathbf{E}\mathbf{q} + \mathbf{r}_{i} = \mathbf{q} + \mathbf{r}_{i}
\tag{3.5.17}
\end{equation*}
因此
\begin{equation*}
\{R_{i} \mid \mathbf{r}_{i}\}^{2}\mathbf{q} = \mathbf{q} + \mathbf{r}_{i} + \mathbf{A}\mathbf{r}_{i}
\tag{3.5.18}
\end{equation*}
且
\begin{equation*}
\begin{aligned}
\{R_{i} \mid \mathbf{r}_{i}\}^{n}\mathbf{q} &= \mathbf{q} + \mathbf{r}_{i} + \mathbf{A}\mathbf{r}_{i} + \mathbf{A}^{2}\mathbf{r}_{i} + \dots + \mathbf{A}^{n-1}\mathbf{r}_{i}\\
&= \mathbf{q} + n\mathbf{w}
\end{aligned}
\tag{3.5.19}
\end{equation*}
因为 $n$ 次螺旋转动操作等价于纯平移 $n\mathbf{w}$。整理式 (3.5.19) 得螺旋矢量 $\mathbf{w}$ 的显式表达式
\begin{equation*}
\mathbf{w} = (1/n)(\mathbf{E} + \mathbf{A} + \mathbf{A}^{2} + \dots + \mathbf{A}^{n-1})\mathbf{r}_{i}.
\tag{3.5.20}
\end{equation*}
螺旋轴的方向当然由 $\mathbf{w}$ 的方向给出，但由式 (3.5.16) 更容易求得该方向。轴的定位则可由“轴上的点 $\mathbf{x}$ 满足 $\mathbf{A}\mathbf{x} + \mathbf{r}_{i} = \mathbf{x} + \mathbf{w}$”这一事实确定，即
\begin{equation*}
(\mathbf{A} - \mathbf{E})\mathbf{x} = \left(\frac{1}{n}\mathbf{B} - \mathbf{E}\right)\mathbf{r}_{i}
\tag{3.5.21}
\end{equation*}
其中
\begin{equation*}
\mathbf{B} = \mathbf{E} + \mathbf{A} + \mathbf{A}^{2} + \dots + \mathbf{A}^{n-1}.
\tag{3.5.22}
\end{equation*}
若 $n = 2$，式 (3.5.21) 有特别简单的解 $\mathbf{x} = \frac{1}{2}\mathbf{r}_{i}$。若 $\det(\mathbf{A}) = -1$，则要考虑三种可能的情形：(a) $I\ ( = IC_{1})$；(b) $\sigma\ ( = IC_{2})$；(c) $S_{6}\ ( = IC_{3})$；$S_{4}\ ( = IC_{4})$；以及 $S_{3}\ ( = IC_{6})$。

\noindent(a) $I\ ( = IC_{1})$。反演中心 $\mathbf{x}$ 可由 $\mathbf{x} = \{I \mid \mathbf{r}_{i}\}\mathbf{x} = -\mathbf{x} + \mathbf{r}_{i}$ 求得，故 $\mathbf{x} = \frac{1}{2}\mathbf{r}_{i}$。

\noindent(b) $\sigma\ ( = IC_{2})$。平面的取向可以用一个过原点且垂直于该平面的矢量 $\mathbf{q}$ 描述；$\mathbf{q}$ 满足 $\mathbf{A}\mathbf{q} = -\mathbf{q}$，即
\begin{equation*}
(\mathbf{A} + \mathbf{E})\mathbf{q} = 0.
\tag{3.5.23}
\end{equation*}
对滑移面，滑移矢量 $\mathbf{w}$ 可由下述事实求得：把 $\{R_{i} \mid \mathbf{r}_{i}\}$ 相继作用两次于平面内的点 $\mathbf{x}$，把它移到 $\mathbf{x} + 2\mathbf{w}$，于是由式 (3.5.18)，
\begin{equation*}
\{R_{i} \mid \mathbf{r}_{i}\}^{2}\mathbf{x} = \mathbf{x} + \mathbf{r}_{i} + \mathbf{A}\mathbf{r}_{i} = \mathbf{x} + 2\mathbf{w}.
\tag{3.5.24}
\end{equation*}
因此
\begin{equation*}
\mathbf{w} = \frac{1}{2}(\mathbf{A} + \mathbf{E})\mathbf{r}_{i}.
\tag{3.5.25}
\end{equation*}
平面内任一点 $\mathbf{x}$ 满足 $\{R_{i} \mid \mathbf{r}_{i}\}\mathbf{x} = \mathbf{A}\mathbf{x} + \mathbf{r}_{i} = \mathbf{x} + \mathbf{w}$，故 $\mathbf{x}$ 满足
\begin{equation*}
(\mathbf{A} - \mathbf{E})\mathbf{x} = \frac{1}{2}(\mathbf{A} - \mathbf{E})\mathbf{r}_{i},
\tag{3.5.26}
\end{equation*}
它显然有特解 $\mathbf{x} = \frac{1}{2}\mathbf{r}_{i}$。

\noindent(c) $S_{6}\ ( = IC_{3})$；$S_{4}\ ( = IC_{4})$；$S_{3}\ ( = IC_{6})$。若 $\mathbf{q}$ 是过原点且平行于其中某根轴的矢量，则 $\mathbf{A}\mathbf{q} = -\mathbf{q}$，$\mathbf{q}$ 同样由求解
\begin{equation*}
(\mathbf{A} + \mathbf{E})\mathbf{q} = 0.
\tag{3.5.23}
\end{equation*}
得到。\footnote{中译者注：原书对情形 (b) 与 (c) 重复使用了编号 (3.5.23)，照录。}不存在螺旋矢量 $\mathbf{w}$。反演中心可由 $\{R_{i} \mid \mathbf{r}_{i}\}\mathbf{x} = \mathbf{A}\mathbf{x} + \mathbf{r}_{i} = \mathbf{x}$ 求得，即
\begin{equation*}
(\mathbf{E} - \mathbf{A})\mathbf{x} = \mathbf{r}_{i}.
\tag{3.5.27}
\end{equation*}

在《国际表》中每个矩阵 $\mathbf{A}$ 都属于表 3.9 所示的四种不同类型之一。确定了 $R_{i}$ 的形式和 $\mathbf{A}$ 的类型之后，就可以利用表 3.10（它用到式 (3.5.16)--(3.5.27)）来确定与《国际表》中该空间群每个等价位置相对应的对称元素的取向和位置。

\begin{table}[H]
\centering
\caption*{\textbf{表 3.9}\ 由《X 射线晶体学国际表》产生的矩阵 $\mathbf{A}$ 的类型}
\begin{tabular}{@{}ll@{}}
\toprule
(i) & $\mathbf{A} = \begin{pmatrix} +1 & 0 & 0\\ 0 & \pm 1 & 0\\ 0 & 0 & \pm 1 \end{pmatrix}$；\qquad
(iii) (a) $\mathbf{A} = \begin{pmatrix} 0 & 0 & \pm 1\\ +1 & 0 & 0\\ 0 & \pm 1 & 0 \end{pmatrix}$，\\[3ex]
(ii) (a) & $\mathbf{A} = \begin{pmatrix} \pm 1 & 0 & 0\\ 0 & 0 & \pm 1\\ 0 & -1 & 0 \end{pmatrix}$，\qquad\quad
(b) $\mathbf{A} = \begin{pmatrix} 0 & \pm 1 & 0\\ 0 & 0 & \pm 1\\ +1 & 0 & 0 \end{pmatrix}$；\\[3ex]
(b) & $\mathbf{A} = \begin{pmatrix} 0 & 0 & \pm 1\\ 0 & \pm 1 & 0\\ \pm 1 & 0 & 0 \end{pmatrix}$，\qquad\quad
(iv) $\mathbf{A} = \begin{pmatrix} a_{11} & a_{12} & 0\\ a_{21} & a_{22} & 0\\ 0 & 0 & \pm 1 \end{pmatrix}$\quad
其中有一个 $a_{ik} = 0$，\\
\multicolumn{2}{l}{其余三个 $a_{ik} = \pm 1$；$i, k = 1, 2$。}\\
\bottomrule
\end{tabular}
\end{table}

\noindent\textbf{表 3.9 的注}\quad（i）假定采用的是 Henry and Lonsdale（1965）给出的定向；对其他定向，将不得不求解式 (3.5.16)--(3.5.27) 中的相应方程。

\begin{figure}[H]
\centering
\includegraphics[width=0.92\textwidth]{c3_t310_p1.png}
\caption*{\textbf{表 3.10}\ 对称元素的测定（据 Wondratschek and Neubüser (1967)，原书第 140 页起扫描；第 1 页）。}
\end{figure}
\begin{figure}[H]
\centering
\includegraphics[width=0.92\textwidth]{c3_t310_p2.png}
\caption*{\textbf{表 3.10}（第 2 页）。}
\end{figure}
\begin{figure}[H]
\centering
\includegraphics[width=0.92\textwidth]{c3_t310_p3.png}
\caption*{\textbf{表 3.10}（第 3 页，页末为表的注）。}
\end{figure}

表 3.10 的注（译文）：

\noindent（i）与《国际表》一致，$a$、$b$、$c$ 是惯用晶胞交于一个顶点的三条棱的长度；$\mathbf{a}$、$\mathbf{b}$、$\mathbf{c}$ 是沿这些棱的矢量。

\noindent（ii）矩阵 $\mathbf{A}(a_{ij})$ 与矢量 $\mathbf{r}_{i}$ 在正文中定义，见式 (3.5.14) 与 (3.5.15)。

\noindent（iii）$\mathbf{A}$ 的类型（i）、（ii）、（iii）或（iv）要用表 3.9 确定。

\noindent（iv）$C_{n}$ 是转角为 $2\pi/n$ 的纯转动。

\noindent（v）本表适用于 Henry and Lonsdale（1965）给出的晶胞定向；对其他定向，将不得不求解式 (3.5.16)--(3.5.27) 中的相应方程。

\noindent\textbf{例 3.5.2}\quad 为了说明表 3.10 的用法，我们用在例 3.5.1 中确定的矩阵 $(\mathbf{A}, \mathbf{r})$：
\begin{equation*}
(\mathbf{A}, \mathbf{r}) = \begin{pmatrix} 1 & 0 & 0 & \frac{1}{2}\\ 0 & -1 & 0 & \frac{1}{2}\\ 0 & 0 & -1 & 0\\ 0 & 0 & 0 & 1 \end{pmatrix}
\tag{3.5.28}
\end{equation*}
于是
\begin{equation*}
\mathbf{A} = \begin{pmatrix} 1 & 0 & 0\\ 0 & -1 & 0\\ 0 & 0 & -1 \end{pmatrix}
\tag{3.5.29}
\end{equation*}
从而
\begin{equation*}
\chi(\mathbf{A}) = -1
\tag{3.5.30}
\end{equation*}
且
\begin{equation*}
\det(\mathbf{A}) = +1.
\tag{3.5.31}
\end{equation*}
由表 3.8，$\{R_{i} \mid \mathbf{r}_{i}\}$ 是一个 $C_{2}$ 型转动；再由表 3.10，由于 $i = 1$，它沿 $x$ 方向，即 $C_{2x}$。螺旋矢量的方向为 $\mathbf{r}_{1}\mathbf{q} = \frac{1}{2}\mathbf{a}$，轴的位置由其通过点 $\left(\frac{1}{2}r_{1}, \frac{1}{2}r_{2}, \frac{1}{2}r_{3}\right)$，即通过 $(\frac{1}{4}a, \frac{1}{4}b, 0)$ 这一事实确定。
